% answers_short.tex —— 2006 年·普通高等学校招生考试·湖北卷文科 · 简短答案速查
% 由 tools/make_short_answers.py 自 answers.tex 的 \gksolhead/\gkanswer 提取，
% 并与 2006/2006_数学真题_湖北文_答案要点.md 逐题比对后生成，勿手改（改源 MD 后重跑）。

\begin{gknotice}
\textbf{使用说明}：本节只给最终答案便于快速核对，完整推导见解析版。答案为编写者独立推导，\textbf{非官方标准答案}。
\end{gknotice}

\gkpart{选择题}
\begin{center}\begin{tabular}{ccccccccccc}
\toprule
题号 & 1 & 2 & 3 & 4 & 5 & 6 & 7 & 8 & 9 & 10 \\
\midrule
答案 & C & D & A & A & B & D & B & C & D & A \\
\bottomrule
\end{tabular}\end{center}

\gkpart{填空题}
\begin{itemize}[leftmargin=2.6em]
  \item[\textbf{11.}] $\dfrac{\sqrt3}{2}$
  \item[\textbf{12.}] U/PENDING；仅按题面给出的边缘概率不能确定联合事件概率。若五人反应相互独立，条件答案为 $0.94$
  \item[\textbf{13.}] $0<k<\dfrac43$
  \item[\textbf{14.}] $78$
  \item[\textbf{15.}] $\left(\dfrac43\pi R^3\right)'=4\pi R^2$；球的体积函数的导数等于球的表面积函数
\end{itemize}

\gkpart{解答题}
\begin{itemize}[leftmargin=2.6em]
  \item[\textbf{16.}] (1) 最大值 $\dfrac{3+\sqrt2}{2}$，最小正周期 $\pi$；(2) $x\in[-\dfrac\pi8+k\pi,\dfrac{3\pi}8+k\pi]$，$k\in\mathbb Z$
  \item[\textbf{17.}] (1) 青年人40\%，中年人50\%，老年人10\%；(2) 青年人60人，中年人75人，老年人15人
  \item[\textbf{18.}] (1) $\dfrac1{\sqrt5}$；(2) $1$
  \item[\textbf{19.}] $a=-c$，$b=-2c-3$；递增、递减区间按 $c>-3$ 或 $c<-3$ 分类，见解析正文。题设排除 $c=-3$
  \item[\textbf{20.}] (1) $a_n=6n-5$；(2) $m=10$
  \item[\textbf{21.}] (1) $\dfrac{x^2}{4}+\dfrac{y^2}{3}=1$；(2) 见解析版
\end{itemize}
